\documentclass[book.tex]{subfiles}

\begin{document}

\chapter{The Transfer Matrix}

\label{chap:transfer_matrix}
\noindent\rule{11cm}{0.4pt} \\
\textbf{Prerequisites:}  \autoref{chap:schrodinger} \\
\textbf{Difficulty Level:} *** \\
\noindent\rule{11cm}{0.4pt}
\vspace{5mm}

\textcolor{red}{TODO: This is mainly notes from Shankar. May or or may not want to turn into full chapter}
Consider a 1-dimensional Ising model
\begin{align}
    Z = \sum_{s_i} \prod_i e^{K(s_i s_{i+1}-1)}
\end{align}
We define the transfer matrix as follows
\begin{align}
    T_{s's} &= e^{K(s's-1)} \\
    T &= \begin{pmatrix}
    1 & e^{-2K} \\ 
    e^{-2K} & 1 \\
    \end{pmatrix}
\end{align}
We can rewrite partition function $Z$ in terms of transfer matrix $T$ as
\begin{align}
    Z &= \sum_{s_1=1,\dotsc,N-1} T_{s_Ns_{N-1}}\dotsc T_{s_2s_2}T_{s_1s_0} \\
    &= \langle s_N | T^N s_0 \rangle
\end{align}
In the case that we have periodic boundary conditions where $s_0 = s_N$, we then obtain
\begin{align}
    Z &= \mathrm{Tr} T^N
\end{align}

\textcolor{red}{TODO: Shankar has a lot of description of the Perron-Frobenius decomposition. Worth covering? Also, should this identity be derived?}

\begin{align}
    T &= e^{K*\sigma_1}
\end{align}
If you add a magnetic field $h$ to the Ising model, then you get
\begin{align}
    T&= e^{K*sigma_1}e^{h\sigma_3} \\
    &= T_K T_h
\end{align} 


\subsection{The Hamiltonian}

The Hamiltonian $H$ is defined as 
below
\begin{align}
    T &= e^{-H}\\
    H &= -K^* \sigma_1
\end{align}

\subsection{Discussion}

\textcolor{red}{TODO: We may want to discuss how the transfer matrix can act as a bridge between the partition function and quantum mechanical systems. I don't fully get this myself yet.}
\end{document}